\documentclass[../main.tex]{subfiles}
\begin{document}
\chapter{2018-2019学年微积分（一）（上）期中考试}

\section{基本计算题(每小题 6 分，共 60 分)}
\begin{enumerate}
    \item 求极限$l=\lim_{n\to \infty}\sum_{k=1}^{n}\frac{1}{n+k^{\alpha}}$，其中$\alpha
              \in(0,1)$.
          \vspace{11em}

    \item 求极限$l=\lim_{x\to 0}\frac{\ln\left(\mathrm{e}^{2\tan x}+\sin x^{2}\right)}{x}$.
          \vspace{10em}

    \item 求极限$l=\lim_{x\to 0^+}\left(\mathrm{e}^{x}-1\right)^{\frac{1}{\ln x}}$.
          \vspace{11em}

    \item 已知$\lim_{x\to1}\frac{ax+bx+b}{\sqrt{3x+1}-\sqrt{x+3}}=2$，求常数$a,b$的值.
          \vspace{10em}

    \item 设可微函数$y=y(x)$由方程$y+y\mathrm{e}^{x}=2\cos y\sin x-4x$确定，求$y'
              (0)$.
          \vspace{10em}

    \item 求曲线$r=2\sin3\theta$在$\theta=\frac{\pi}{3}$处的切线方程.
          \vspace{10em}

    \item 设函数$f(x)=x\sqrt{1-x^{2}}+\arcsin x$，求微分$\left.\mathrm{d}y\right|
              _{x=0}$.
          \vspace{11em}

    \item 设函数$y=x+x^{3}$的反函数为$x=g(y)$，求$g''(2)$.
          \vspace{12em}

    \item 设$y=x^{2}\cos 2x$，求$y^{(5)}\left(\frac{\pi}{2}\right)$.
          \vspace{8em}

    \item 设函数$y=y(x)$由方程$\begin{cases}
                  x=t\cos t-\sin t, \\
                  y=\cos t
              \end{cases}$确定，求$\frac{\mathrm{d}y}{\mathrm{d}x}$和$\frac{\mathrm{d}^{2}y}{\mathrm{d}x^{2}}$.
          \vspace{8em}
\end{enumerate}

\section{综合题(每小题 6 分，共 30 分)}
\begin{enumerate}
    \setcounter{enumi}{10}
    \item 讨论函数$f(x)=\frac{\sin\pi x}{x^{2}-1}$的连续性，并对函数的间断点判别类型.
          \vspace{8em}

    \item 求无穷小量$u(x)=\left(\frac{1+2\cos x}{3}\right)^{x^2}-1(x\to0)$的主部和阶数.
          \vspace{8em}

    \item 求函数与$f(x)=
              \begin{cases}
                  \ln(1+x), & x\geq0, \\
                  x,        & x<0
              \end{cases}$的导函数$f'(x)$，并讨论$f'(x)$的连续性.
          \vspace{9em}

    \item 已知$a>1$，$n\geq1$，证明不等式$\frac{\sqrt[n]{a}-\sqrt[n+1]{a}}{\ln
                  a}<\frac{\sqrt[n]{a}}{n^{2}}$.
          \vspace{10em}

    \item 一架飞机在$H$米高空以$a$米/秒的速度水平匀速飞行.设在$t=0$时刻有一探照灯位于飞机正下方的地面上跟踪飞机.问$t$秒以后探照灯应以怎样的角速度转动才能照到飞机？
          \vspace{12em}
\end{enumerate}

\section{证明题(每小题 5 分，共 10 分)}
\begin{enumerate}
    \setcounter{enumi}{15}
    \item 设数列$\{x_{n}\}$由递推公式$x_{1}=1$，$x_{n+1}=x_{n}+\frac{2}{x_{n}^{2}}$给出.证明数列$\{
              x_{n}\}$无界.
          \vspace{12em}

    \item 设$f(x)$在区间$[a,b]$有定义并且对于任何$x,y\in[a,b](x\neq y)$成立
          \[
              |f(x)-f(y)|\leq M(x-y)^{2},
          \]

          其中$M$为常数.证明$f(x)$在$[a,b]$上恒为常数.
\end{enumerate}

\chapter{2019-2020学年微积分（一）（上）期中考试参考答案}
\section{基本计算题(每小题 6 分，共 60 分)}
\begin{enumerate}
    \item \textbf{Solution}. 因$\frac{n}{n+n^{\alpha}}<\sum_{k=1}^{n}\frac{1}{n+k^{\alpha}}
              <1$，而$\lim_{n\to\infty}\frac{n}{n+n^{\alpha}}=\lim_{n\to\infty}\frac{1}{1+n^{\alpha-1}}
              =1$，

          由夹逼定理知$l=1$.

    \item \textbf{Solution}.
          \[
              \begin{aligned}
                  l=\lim_{x\to0}\frac{\ln\left(\mathrm{e}^{2\tan x}+\sin x^2\right)}{x}= & \lim_{x\to0}\frac{\mathrm{e}^{2\tan x}+\sin x^2-1}{x}        \\
                  =                                                                      & \lim_{x\to0}\frac{2\tan x+\sin x^2}{x}                       \\
                  =                                                                      & \lim_{x\to0}\frac{2\tan x}{x}+\lim_{x\to0}\frac{\sin x^2}{x} \\
                  =                                                                      & 2.
              \end{aligned}
          \]

    \item \textbf{Solution}.
          \[
              \begin{aligned}
                  l= & \lim_{x\to0^+}\left(\mathrm{e}^{x}-1\right)^{\frac{1}{\ln x}}= \mathrm{e}^{\lim\limits_{x\to0^+}\frac{\ln\left(\mathrm{e}^x-1\right)}{\ln x}}        \\
                  =  & \mathrm{e}^{\lim\limits_{x\to0^+}\frac{\frac{\mathrm{e}^x}{\mathrm{e}^x-1}}{\frac{1}{x}}}= \mathrm{e}^{\lim\limits_{x\to0^+}\frac{x\mathrm{e}^x}{x}} \\
                  =  & \mathrm{e}^{\lim\limits_{x\to0^+}\mathrm{e}^x}= \mathrm{e}.
              \end{aligned}
          \]

    \item \textbf{Solution}.计算
          \[
              \begin{aligned}
                  \lim_{x\to1}\frac{ax+bx+b}{\sqrt{3x+1}-\sqrt{x+3}}= & \lim_{x\to1}\frac{[(a+b)x+b]\left(\sqrt{3x+1}+\sqrt{x+3}\right)}{(3x+1)-(x+3)} \\
                  =                                                   & \lim_{x\to1}\frac{2[(a+b)x+b]}{x-1}                                            \\
                  =                                                   & 2.
              \end{aligned}
          \]

          因此有$\begin{cases}
                  a+2b=0, \\
                  a+b=1.
              \end{cases}$ 解得$a=2,b=-1$.

    \item \textbf{Solution}. 方程两边对$x$求导得到
          \[
              y'+y'\mathrm{e}^{x}+y\mathrm{e}^{x}=-2\sin y\sin x\cdot y'+2\cos y\cos x-
              4,
          \]

          代入$x=0,y=0$解得$y'(0)=-1$.

    \item \textbf{Solution}. 曲线的参数方程为
          \[
              \begin{cases}
                  x=2\sin3\theta\cos\theta, \\
                  y=2\sin3\theta\sin\theta.
              \end{cases}
          \]

          所以$\frac{\mathrm{d}y}{\mathrm{d}x}=\frac{\frac{\mathrm{d}y}{\mathrm{d}\theta}}{\frac{\mathrm{d}x}{\mathrm{d}\theta}}
              =\frac{6\cos3\theta\sin\theta+2\sin3\theta\cos\theta}{6\cos3\theta\cos\theta-2\sin3\theta\sin\theta}$.

          代入$\theta=\frac{\pi}{3}$得切线斜率$k=\sqrt{3}$，又切点坐标为$(0,0)$，所以切线方程为$y
              =\sqrt{3}x$.

    \item \textbf{Solution}. $f'(x)=\sqrt{1-x^{2}}+\frac{-2x^{2}}{2\sqrt{1-x^{2}}}
              +\frac{1}{\sqrt{1-x^{2}}}=2\sqrt{1-x^{2}}$，

          所以$\left.\mathrm{d}y\right|_{x=0}=f'(0)\,\mathrm{d}x=2\,\mathrm{d}x$.

    \item \textbf{Solution}. $g'(y)=\frac{\mathrm{d}x}{\mathrm{d}y}=\frac{1}{\frac{\mathrm{d}y}{\mathrm{d}x}}
              =\frac{1}{1+3x^{2}}$，

          $g''(y)=\frac{\mathrm{d}}{\mathrm{d}y}\left(g'(y)\right)=\frac{\mathrm{d}}{\mathrm{d}x}
              \left(g'(y)\right)\cdot\frac{\mathrm{d}x}{\mathrm{d}y}=-\frac{6x}{(1+3x^{2})^{3}}$.

          代入$x=1,y=2$得$g''(2)=-\frac{3}{32}$.

    \item \textbf{Solution}. 利用Leibniz法则得到
          \[
              \begin{aligned}
                  y^{(5)}= & x^{2}\left(\cos 2x\right)^{(5)}+5\cdot 2x\left(\cos 2x\right)^{(4)}+10\cdot 2\left(\cos 2x\right)^{(3)} \\
                  =        & -32x^{2}\sin 2x+160x\cos 2x+160\sin 2x.
              \end{aligned}
          \]

          代入$x=\frac{\pi}{2}$得$y^{(5)}\left(\frac{\pi}{2}\right)=-80\pi$.

    \item \textbf{Solution}. $\frac{\mathrm{d}y}{\mathrm{d}x}=\frac{\frac{\mathrm{d}y}{\mathrm{d}t}}{\frac{\mathrm{d}x}{\mathrm{d}t}}
              =\frac{-\sin t}{\cos t-t\sin t-\cos t}=\frac{-\sin t}{-t\sin t}=\frac{1}{t}$.

          $\frac{\mathrm{d}^{2}y}{\mathrm{d}x^{2}}=\frac{\mathrm{d}}{\mathrm{d}t}\left
              (\frac{\mathrm{d}y}{\mathrm{d}x}\right)\cdot\frac{\mathrm{d}t}{\mathrm{d}x}
              =-\frac{1}{t^{2}}\cdot\frac{1}{-t\sin t}=\frac{1}{t^{3}\sin t}$.
\end{enumerate}

\section{综合题(每小题 6 分，共 30 分)}
\begin{enumerate}
    \setcounter{enumi}{10}
    \item \textbf{Solution}. 函数$f(x)$的间断点为$-1,1$.因
          \[
              \begin{aligned}
                  \lim_{x\to-1}\frac{\sin\pi x}{x^2-1}= & \lim_{x\to-1}\frac{\pi\cos\pi x}{2x}=\frac{\pi\cos(-\pi)}{-2}=\frac{\pi}{2}, \\
                  \lim_{x\to1}\frac{\sin\pi x}{x^2-1}=  & \lim_{x\to1}\frac{\pi\cos\pi x}{2x}=\frac{\pi\cos\pi}{2}=-\frac{\pi}{2}.
              \end{aligned}
          \]

          所以$x=-1,1$均为可去间断点.

    \item \textbf{Solution}.
          \[
              \begin{aligned}
                  u=\left(\frac{1+2\cos x}{3}\right)^{x^2}-1= & \exp\left({x^2\ln\frac{1+2\cos x}{3}}\right)-1 \\
                  \sim                                        & x^{2}\ln\frac{1+2\cos x}{3}                    \\
                  \sim                                        & \frac{2x^2}{3}\left(\cos x-1\right)            \\
                  \sim                                        & -\frac{1}{3}x^{4}.
              \end{aligned}
          \]

          所以$u$的主部为$-\frac{1}{3}x^{4}$，阶数为4.

    \item \textbf{Solution}. 当$x>0$时，$f'(x)=\frac{1}{1+x}$；当$x<0$时，$f
              '(x)=1$；

          当$x=0$时，$\lim_{x\to0^+}f'(x)=1=\lim_{x\to0^-}f'(x)$，所以$f'(0)=1$.

          因此
          \[
              f'(x)=
              \begin{cases}
                  \frac{1}{1+x}, & x>0,    \\
                  1,             & x\leq0.
              \end{cases}
          \]

          当$x\neq0$时，$f'(x)$显然连续；

          当$x=0$时，$\lim_{x\to0^+}f'(x)=1=f'(0)=\lim_{x\to0^-}f'(x)$，所以$f'(x)$在$x
              =0$处连续.

          因此$f'(x)$在$\mathbf{R}$上处处连续.

    \item \textbf{Solution}. 令$f(x)=a^{x}$，显然$f(x)$在区间$\left[\frac{1}{n+1}
                  ,\frac{1}{n}\right]$内可导.由Lagrange中值定理知存在$\xi\in\left(\frac{1}{n+1}
              ,\frac{1}{n}\right)$使得
          \[
              \sqrt[n]{a}-\sqrt[n+1]{a}=f\left(\frac{1}{n}\right)-f\left(\frac{1}{n+1}\right
              )=f'(\xi)\left(\frac{1}{n}-\frac{1}{n+1}\right)=\frac{a^{\xi}\ln a}{n(n+1)}
              <\frac{a^{\frac{1}{n}}\ln a}{n^{2}},
          \]

          即$\frac{\sqrt[n]{a}-\sqrt[n+1]{a}}{\ln a}<\frac{\sqrt[n]{a}}{n^{2}}$.

    \item \textbf{Solution}. 设飞机飞过的距离为$x$米，飞机与探照灯的连线与竖直方向的夹角为$\theta$.

          由几何关系有$x=H\tan\theta$，故
          \[
              \frac{\mathrm{d}x}{\mathrm{d}t}=a=H\sec^{2}\theta\cdot \frac{\mathrm{d}\theta}{\mathrm{d}t}
              .
          \]

          当前时刻$x=at$，故$\sec\theta=\frac{\sqrt{H^{2}+a^{2}t^{2}}}{H}$，代入上式得
          \[
              \frac{\mathrm{d}\theta}{\mathrm{d}t}=\frac{aH}{H^{2}+a^{2}t^{2}}.
          \]
\end{enumerate}

\section{证明题(每小题 5 分，共 10 分)}
\begin{enumerate}
    \setcounter{enumi}{15}
    \item \textbf{Proof}.

          显然$x_{n+1}-x_{n}=\frac{2}{x_{n}^{2}}>0$，即$\{x_{n}\}$严格单调增加.

          若$\{x_{n}\}$有界，则$\{x_{n}\}$收敛.设$\lim_{n\to\infty}x_{n}=a$，则$a=a+\frac{2}{a^{2}}$，此方程无解，矛盾.

          因此$\{x_{n}\}$无界.

    \item \textbf{Proof}. 由题可知，$\forall\,x,y\in[a,b](x\neq y)$，有$0\leq
              \left|\frac{f(x)-f(y)}{x-y}\right|\leq M|x-y|$.

          当$x\to y$时，$M|x-y|\to0$，由夹逼定理知$\lim_{x\to y}\left|\frac{f(x)-f(y)}{x-y}
              \right|=0$，即$f'(x)=0$.

          所以$f(x)$在$[a,b]$上恒为常数.
\end{enumerate}
\end{document}




